\documentclass{article}
\usepackage{amsmath,hyperref}
\title{Finishing the first crossing releases the gate for a second job}
\author{}
\date{}
\begin{document}
\maketitle
\paragraph{Joint held-in scenario.}
The
\href{https://github.com/eternal-november/eos-research-programmes/pull/806}{expanded-body cave replay}
already returned a scoped report and produced the corresponding
checked egress route for courier Q. The present development test
adds a second owed crossing W through the gate which that reply
classifies as closed. The actuator resource has two mutually exclusive
open modes (A or B); it may switch only after Q's last contact with
the first open gate, followed by an optional all-closed interval
\(\gamma\). This is a declared interlock and clock, not actual door
dynamics or a protocol for granting the interlock. The separate
inspection report remains paid under the native parent execution.
W's end is its required local service point, not an additional
whole-cave egress.

The old and new trajectories are held in one mode history, not
evaluated under independently favorable panel states. W follows a
5 m straight line from \((99,66,8)\) to \((99,61,8)\) when it needs
A, or from \((65,53,6)\) to \((60,53,6)\) when it needs B.
The exact triangle tests prove positive clearance against the pinned
placements and stationary bodies for its actual open panel; a
strict interior witness proves it needs that panel open. W's
entire expanded swept box is disjoint from every expanded swept
box of Q's chosen route, so their motions can overlap outside
the gate-mode restriction. Dispatch into a confirmed open mode,
not re-use of an old report alone, licenses W's crossing.

\paragraph{Earliest service on the declared paths.}
Let \(q\) be Q's last contact, computed with the enlarged box by
the parent's exact slab and acceleration calculation. The first
permitted panel switch on the model clock is the first sixteenth
of a second strictly after \(q\); W's panel opens \(\gamma\) later.
For W let \(L=5\) and \(x\) be progress along its required segment.
From rest-endpoint conditions, monotone motion and unit acceleration
magnitude, its speed is at most
\(\min(\sqrt{2x},3,\sqrt{2(L-x)})\).
The standard full accelerate/brake trajectory attains this speed
profile. If \(x_*\) is first contact with W's gate box, the integral
from \(x_*\) to \(L\) of one over this speed is a lower bound
on the time needed after the mode opens. Starting the same trajectory
early so it reaches first contact exactly at the opening attains
that bound. Thus the gate allocation rule need not make W stand
still until the signal, although without verified permission
it cannot cross ahead of the opening. The ordering convention is
actuate the switch first at the appointed cut, then admit contact;
an earlier sample remains outside the closed box.

At \(\gamma=1\) in the compact parent arm, the measured
source-derived schedule puts W at its service point at
55.13711535 s (A initially closed) or 31.54798979 s (B
initially closed). Q still reaches staging under 200 s.
Keeping the mode through all of Q's egress would delay the
same W service by about 67.31 or 113.38 s. The full-record
parent arm also supports both services, one second later in
the compared scheduling. These are observed-model service
frontiers, not acceptance deadlines chosen to force a win.

\paragraph{Before the crossing and after it are not the same delay.}
The admitted delay probe inserts a wait at an existing
rest waypoint: before the unique gate-contact leg or just
after it. Both can delay Q's ultimate completion equally.
Only the former shifts the latest gate overlap and the earliest
safe appointment for W. Reusing the nominal flip slot under
any tested positive upstream delay fails to preserve the gate
requirement; an adverse switch can occur before or during the
shifted obligation. Delays after the last contact do not consume
the gate again. Ninety-six combinations cover the two native
reply schedules, two layouts, three interlock times and
both placements of four dwell sizes. A missing report does
not by itself cancel Q's pre-existing crossing authorization.

\paragraph{Comparator and limits.}
A direct supervisor with the same certified last-contact bound,
mode authority and eligible clock realizes the same schedule.
Ordinary action-dependency execution can release on actual
clearance rather than the final egress; see
\href{https://whoenig.github.io/publications/2019_RA-L_Hoenig.pdf}{H\"onig et al., section IV-B}.
The whole-egress comparator demonstrates excess grant scope,
not an advantage over that stronger mechanism. This test uses
a prearranged common modeled clock and timing envelope and
does not price initial dissemination of that appointment.
If a body stops unexpectedly in the gate or can start late,
the bounded certificate no longer suffices. Independent
selection, transport latency, physical actuator confirmation
and global route optimality remain unverified.

Run \texttt{python3 gate\_queue.py} with the pinned parent input
directories or in the complete offline bundle.
\end{document}
